% !TEX program = xelatex
% 공통수학2 · Ⅱ. 집합과 명제 · 2. 명제 · 03 명제의 증명 · 2/2차시
% 학생용 활동지 (답안 없음)
\documentclass[11pt,a4paper]{article}
\usepackage{kotex}
\usepackage{iftex}
\ifPDFTeX\else
  \setmainhangulfont{Noto Serif CJK KR}
  \setsanshangulfont{Noto Sans CJK KR}
\fi
\usepackage[margin=2.1cm]{geometry}
\usepackage{parskip}
\usepackage{enumitem}
\usepackage{array}
\usepackage{tabularx}
\usepackage{xcolor}
\usepackage{amssymb}
\usepackage{tikz}
\usepackage[most]{tcolorbox}
\usepackage{fancyhdr}
\usepackage{titlesec}

\definecolor{goldc}{HTML}{B07C22}
\definecolor{goldstrong}{HTML}{8A5F16}
\definecolor{indigoc}{HTML}{2B3B57}
\definecolor{indigosoft}{HTML}{6E82A6}
\definecolor{linec}{HTML}{D6DACB}
\definecolor{surface2}{HTML}{E4E8DC}
\definecolor{writeline}{HTML}{C7CCBA}
\definecolor{inksoft}{HTML}{52604F}

\titleformat{\section}{\normalfont\sffamily\bfseries\large\color{indigoc}}{\thesection}{0.6em}{}
\titlespacing{\section}{0pt}{14pt}{6pt}
\newtcolorbox{goalbox}{enhanced, breakable, colback=white, colframe=goldc, boxrule=0.5pt, leftrule=3pt, arc=2pt, left=10pt, right=10pt, top=8pt, bottom=8pt}
\newtcolorbox{sheetbox}{enhanced, breakable, colback=white, colframe=linec, boxrule=0.6pt, arc=3pt, left=11pt, right=11pt, top=9pt, bottom=9pt}
\newcommand{\writespace}[1]{%
  \par\nobreak\vspace{3pt}
  \foreach \n in {1,...,#1} {%
    \noindent\textcolor{writeline}{\rule{\linewidth}{0.4pt}}\par\vspace{0.62cm}%
  }
}
\newcommand{\fillblank}[1]{\makebox[#1]{\hrulefill}}
\pagestyle{fancy}\fancyhf{}\renewcommand{\headrulewidth}{0pt}
\fancyfoot[L]{\footnotesize\color{inksoft} 공통수학2 · 2026학년도 2학기 · 지평선고등학교}
\fancyfoot[R]{\footnotesize\color{inksoft} 다음 시간 --- 중단원·대단원 마무리(Ⅱ)}

\begin{document}

\noindent
\begin{minipage}[b]{0.62\linewidth}
  {\sffamily\footnotesize\color{indigosoft} 공통수학2 · Ⅱ. 집합과 명제 · 2. 명제}\\[2pt]
  {\sffamily\bfseries\Large 03 명제의 증명 \textperiodcentered\ 31차시}
\end{minipage}\hfill
\begin{minipage}[b]{0.36\linewidth}
  \raggedleft\small\color{inksoft}
  반\ \fillblank{1.6cm} \quad 번호\ \fillblank{1.6cm}\\[4pt]
  이름\ \fillblank{3.6cm}
\end{minipage}

\vspace{4pt}
\noindent{\color{goldc}\rule{\linewidth}{2.2pt}}
\vspace{10pt}

\begin{goalbox}
\textbf{오늘의 목표}\\
절대부등식의 뜻을 알고, 완전제곱식과 산술·기하평균의 관계를 이용하여 간단한 절대부등식을 증명할 수 있다.
\vspace{4pt}
\small\color{inksoft}
$\square$ 자신 있음 \quad $\square$ 조금 헷갈림 \quad $\square$ 처음 봄 \hfill (수업 전 나의 상태에 표시)
\end{goalbox}

\section{생각 열기 --- 항상 성립할까}

\begin{sheetbox}
모든 실수 $x$에 대하여 항상 성립하는 부등식을 다음에서 모두 찾아 보자.

\begin{enumerate}[leftmargin=1.4em, itemsep=2pt]
  \item $x^2+1>0$
  \item $x^2-1>0$
  \item $|x-1|\ge0$
\end{enumerate}
\writespace{3}
\end{sheetbox}

\section{개념 정리 --- 절대부등식}

\begin{sheetbox}
주어진 집합의 모든 원소에 대하여 성립하는 부등식을 \fillblank{3.4cm}이라고 한다.

\vspace{6pt}
\textbf{기본 성질} ($a$, $b$는 실수)\\
① $a>b \Longleftrightarrow a-b\ \fillblank{1.2cm}0$ \qquad ② $a^2\ \fillblank{1.2cm}\ 0$, \; $a^2+b^2\ \fillblank{1.2cm}\ 0$\\
③ $a^2+b^2=0 \Longleftrightarrow a=\ \fillblank{0.9cm}=0$

\vspace{8pt}
\textbf{예제} --- $a$, $b$가 실수일 때, $a^2+b^2\ge ab$가 성립함을 증명하시오.\\
증명: $a^2-ab+b^2=\left(a-\dfrac b2\right)^2+\dfrac34b^2$이고, $\left(a-\dfrac b2\right)^2\ge0$, $\dfrac34b^2\ge0$이므로 $a^2-ab+b^2\ge0$. 따라서 $a^2+b^2\ge ab$. (등호는 $a=b=0$일 때)

\vspace{8pt}
\textbf{산술평균과 기하평균} --- $a>0$, $b>0$일 때, $\dfrac{a+b}{2}\ \fillblank{1.2cm}\ \sqrt{ab}$가 성립한다. (등호는 $a=b$일 때)
\end{sheetbox}

\section{연습 문제}

\begin{sheetbox}
\textbf{1.} $a$, $b$가 실수일 때, 다음 부등식이 성립함을 증명하시오.\\
\hspace*{1.6em}⑴ $(a+b)^2\ge4ab$\\
\hspace*{1.6em}⑵ $a^2-4ab+5b^2\ge0$
\writespace{7}

\vspace{6pt}
\textbf{2.} $a>0$일 때, 다음 부등식이 성립함을 증명하시오.\\
\hspace*{1.6em}$a+\dfrac1a\ge2$
\writespace{5}
\end{sheetbox}

\section{사고의 가시화 --- See · Think · Wonder}

\noindent{\footnotesize\color{inksoft} 오늘 배운 `절대부등식'을 떠올리며 아래 세 칸을 채워 보자.}\\[6pt]
\begin{tabularx}{\linewidth}{@{}XXX@{}}
\textbf{\footnotesize\color{indigoc} See --- 무엇을 보았나?} &
\textbf{\footnotesize\color{indigoc} Think --- 무엇을 생각했나?} &
\textbf{\footnotesize\color{indigoc} Wonder --- 무엇이 궁금한가?} \\[4pt]
\end{tabularx}
\noindent
\begin{minipage}[t]{0.32\linewidth}\writespace{3}\end{minipage}\hfill
\begin{minipage}[t]{0.32\linewidth}\writespace{3}\end{minipage}\hfill
\begin{minipage}[t]{0.32\linewidth}\writespace{3}\end{minipage}

\section{마무리 성찰}

\begin{sheetbox}
\begin{minipage}[t]{0.48\linewidth}
{\footnotesize\color{inksoft} `명제의 증명' 소단원과 대단원 Ⅱ를 마치며 한 문장으로}
\writespace{2}
\end{minipage}\hfill
\begin{minipage}[t]{0.48\linewidth}
{\footnotesize\color{inksoft} 감사 일기 / 유념 한 줄}
\writespace{2}
\end{minipage}
\end{sheetbox}

\end{document}
